\documentclass[11pt, a4paper]{article}

%---------------------------------------------------------
% Lingua, Codifica e Font
%---------------------------------------------------------
\usepackage[italian]{babel}
\usepackage[T1]{fontenc}
\usepackage{lmodern}
\usepackage[utf8]{inputenc}
\usepackage{csquotes}

%---------------------------------------------------------
% Layout Di Pagina E Tipografia
%---------------------------------------------------------
\usepackage[margin=1.5cm, top=1.5cm, bottom=1.5cm]{geometry}
\usepackage{parskip}
\usepackage[verbose=silent]{microtype}
\usepackage{soul}
\usepackage{fancybox}

%---------------------------------------------------------
% Gestione Dei Colori
%---------------------------------------------------------
\usepackage[table, svgnames]{xcolor}
\definecolor{darkgreen}{RGB}{0,100,0}
\definecolor{theoryblue}{RGB}{30,144,255}

%---------------------------------------------------------
% Matematica E Simboli (Fondamentali per ML)
%---------------------------------------------------------
\usepackage{amsmath, amssymb, amsfonts, amsthm}
\usepackage{mathtools}
\usepackage{siunitx}

% BB OTTIMIZZAZIONE MATRICI LUNGHE PER LA STAMPA
\setlength{\arraycolsep}{10pt}
\usepackage{bm}
\usepackage{arydshln}
\usepackage{mathrsfs}
\usepackage{cancel}

% Operatori comuni nel Machine Learning
\DeclareMathOperator*{\argmax}{arg\,max}
\DeclareMathOperator*{\argmin}{arg\,min}
\DeclareMathOperator{\E}{\mathbb{E}}
\DeclareMathOperator{\Var}{Var}
\DeclareMathOperator{\Cov}{Cov}
% Loss
\newcommand{\Loss}{\mathcal{L}}
% Dataset
\newcommand{\Dataset}{\mathcal{D}}
% Norma
\newcommand{\Norma}[1]{\left\lVert#1\right\rVert}

%---------------------------------------------------------
% Grafica, Immagini E Tabelle
%---------------------------------------------------------
\usepackage{graphicx}
\graphicspath{{./Images/}}
\usepackage{float}
\usepackage{booktabs}
\usepackage{array}
\usepackage{multirow}
\usepackage{tabularx}
\usepackage{caption}
\usepackage{subcaption}

% TikZ e PGFPlots
\usepackage{tikz}
\usetikzlibrary{arrows.meta, patterns, angles, quotes, calc, positioning}
\usepackage{pgfplots}
\pgfplotsset{compat=1.18}

%---------------------------------------------------------
% Link, Riferimenti Ed Elenchi
%---------------------------------------------------------
\usepackage[unicode]{hyperref}
\hypersetup{
  linkcolor=blue,
  urlcolor=blue,
  citecolor=blue,
  pdftitle={Machine Learning \& Data Analysis},
  pdfauthor={Francesco Lo Verde},
}
\usepackage{bookmark}
\usepackage[italian, capitalize]{cleveref}
\usepackage{enumitem}
\setlist{noitemsep}

%---------------------------------------------------------
% Personalizzazione Titoli
%---------------------------------------------------------
\usepackage{titlesec}
\titleformat{\section}{\normalfont\LARGE\bfseries}{\thesection}{1em}{}
\titleformat{\subsection}{\normalfont\Large\bfseries}{\thesubsection}{1em}{}

%---------------------------------------------------------
% Algoritmi
%---------------------------------------------------------
\usepackage[ruled,vlined,linesnumbered,italiano,onelanguage]{algorithm2e}

% Comandi custom per algorithm2e
\SetKwInput{Dati}{Input}
\SetKwInput{Risultato}{Output}
\SetKw{Inizializza}{Inizializza}

% Definizione della struttura Se...Allora...
\SetKwIF{Se}{AltrimentiSe}{Altrimenti}{se}{allora}{altrimenti se}{altrimenti}{finese}

% Definizione della struttura Funzione
\SetKwProg{Fn}{Funzione}{:}{fine}

% Definizione dei cicli (AGGIUNTI ORA)
\SetKwFor{Per}{per}{fai}{fine per}
\SetKwRepeat{Ripeti}{ripeti}{finché}

% Gestione commenti
\SetKwComment{Commento}{// }{}

%---------------------------------------------------------
% Inizio Documento
%---------------------------------------------------------
\title{RISOLUZIONE DI UGUAGLIANZE MATRICIALI LINEARI TRAMITE VETTORIZZAZIONE}
\author{FRANCESCO FERRANTE}
\date{}

\begin{document}

\maketitle

\section{Problema e motivazione}
Mentre la risoluzione di sistemi di uguaglianze lineari è piuttosto standard, la soluzione di equazioni matriciali lineari non è così comune. Tuttavia, molti problemi di interesse nel campo del controllo richiedono di essere in grado di gestire equazioni matriciali lineari (LME). In generale, una LME si scrive come:


\begin{equation}
AXC + HX + XG = B
\end{equation}
dove $A \in \mathbb{R}^{n \times m}$, $B \in \mathbb{R}^{n \times q}$, $C \in \mathbb{R}^{q \times q}$, $H \in \mathbb{R}^{n \times m}$, e $G \in \mathbb{R}^{q \times q}$ sono dati e $X \in \mathbb{R}^{m \times q}$ è l'incognita.

\section{Strumenti}
La difficoltà principale nel risolvere direttamente la (1) è fondamentalmente legata al fatto che il prodotto tra matrici non commuta in generale. Pertanto, è generalmente impossibile riscrivere la (1) come $A'X = B'$ per alcune matrici $A'$ e $B'$. D'altra parte, l'algebra lineare ci insegna che le matrici reali non sono altro che vettori euclidei se si guarda bene! Più formalmente, poiché lo spazio vettoriale $\mathbb{R}^{n \times m}$ è uno spazio vettoriale reale a dimensione finita, può essere mappato nello spazio $\mathbb{R}^{nm}$ tramite un isomorfismo, ovvero una trasformazione lineare invertibile. In particolare, tale trasformazione può essere assunta come l'operazione di "vettorizzazione" definita come segue per ogni $A \in \mathbb{R}^{n \times m}$:
\[
\text{vec}(A) =
\begin{bmatrix}
a_{1,1} \\
\vdots \\
a_{n,1} \\
\vdots \\
a_{1,2} \\
\vdots \\
a_{n,2} \\
\vdots \\
a_{1,m} \\
\vdots \\
a_{n,m}
\end{bmatrix}
\]
Ad esempio, se
\[
A = \begin{bmatrix} 1 & 2 & 3 \\ 0 & 4 & 5 \end{bmatrix} \in \mathbb{R}^{2 \times 3}
\]
allora
\[
\text{vec}(A) = \begin{bmatrix} 1 \\ 0 \\ 2 \\ 4 \\ 3 \\ 5 \end{bmatrix} \in \mathbb{R}^6
\]
Chiaramente, $\text{vec}$ soddisfa tutte le belle proprietà di un isomorfismo, cioè è lineare e invertibile. OK, allora vettorizziamo le cose!
\begin{equation}
\text{vec}(AXC + HX + XG) = \text{vec}(B)
\end{equation}
possiamo certamente farlo perché $\text{vec}$ è invertibile. Inoltre, usando la linearità:
\begin{equation}
\text{vec}(AXC) + \text{vec}(HX) + \text{vec}(XG) = \text{vec}(B)
\end{equation}
Ora sappiamo che possiamo vettorizzare le cose, ma non siamo ancora in grado di gestire la (3) a causa degli elementi moltiplicati. Quindi, per farcela, dobbiamo capire come si comporta la vettorizzazione in termini di moltiplicazione tra matrici. A tal fine, dobbiamo introdurre un'altra operazione, che è il prodotto di Kronecker.

Dati $A \in \mathbb{R}^{m \times n}$ e $B \in \mathbb{R}^{p \times q}$, il prodotto di Kronecker $A \otimes B$ è la seguente matrice a blocchi $pm \times qn$:
\[
A \otimes B =
\begin{bmatrix}
a_{11}B & \dots & a_{1n}B \\
\vdots & \ddots & \vdots \\
a_{m1}B & \dots & a_{mn}B
\end{bmatrix}
\]
Ad esempio, se
\[
A = \begin{bmatrix} 1 & 2 \\ 3 & 4 \end{bmatrix}, \quad B = \begin{bmatrix} 0 & 5 \\ 6 & 7 \end{bmatrix}
\]
Allora
\[
A \otimes B =
\begin{bmatrix}
1 \begin{bmatrix} 0 & 5 \\ 6 & 7 \end{bmatrix} & 2 \begin{bmatrix} 0 & 5 \\ 6 & 7 \end{bmatrix} \\
3 \begin{bmatrix} 0 & 5 \\ 6 & 7 \end{bmatrix} & 4 \begin{bmatrix} 0 & 5 \\ 6 & 7 \end{bmatrix}
\end{bmatrix} =
\begin{bmatrix}
0 & 5 & 0 & 10 \\
6 & 7 & 12 & 14 \\
0 & 15 & 0 & 20 \\
18 & 21 & 24 & 28
\end{bmatrix} \in \mathbb{R}^{4 \times 4}
\]
Incredibilmente, l'aver definito questo strano prodotto permette di gestire la vettorizzazione delle moltiplicazioni standard tra matrici. In particolare, le seguenti due proprietà risolvono il problema.

Siano $A \in \mathbb{R}^{k \times l}$, $B \in \mathbb{R}^{l \times m}$, e $C \in \mathbb{R}^{m \times n}$, allora valgono le seguenti identità:
\begin{align}
\text{vec}(ABC) &= (C^\top \otimes A)\text{vec}(B) \tag{4a} \\
\text{vec}(AB) &= (I_m \otimes A)\text{vec}(B) \tag{4b} \\
\text{vec}(AB) &= (B^\top \otimes I_k)\text{vec}(A) \tag{4c}
\end{align}

\section{Soluzione tramite la teoria dei sistemi lineari}
Ora abbiamo tutto ciò di cui abbiamo bisogno per affrontare la (1). In particolare, torniamo alla (3) e vediamo come la (4) può essere utilizzata per liberare l'incognita $X$. Gestiamo termine per termine:
\begin{align*}
\text{vec}(AXC) &\overset{(4a)}{=} (C^\top \otimes A)\text{vec}(X) \\
\text{vec}(HX) &\overset{(4b)}{=} (I_q \otimes H)\text{vec}(X) \\
\text{vec}(XG) &\overset{(4c)}{=} (G^\top \otimes I_m)\text{vec}(X)
\end{align*}
Quindi, mettendo tutto insieme otteniamo:
\begin{equation}
(C^\top \otimes A)\text{vec}(X) + (I_q \otimes H)\text{vec}(X) + (G^\top \otimes I_m)\text{vec}(X) = \text{vec}(B)
\end{equation}
che può essere convenientemente riscritta come:
\[
\left( (C^\top \otimes A) + (I_q \otimes H) + (G^\top \otimes I_m) \right) \text{vec}(X) = \text{vec}(B)
\]
Ma questo è un normale sistema di equazioni lineari!! Quindi abbiamo finito.

\section{Applicazione alle equazioni del regolatore}
Applichiamo ora tutto l'apparato presentato al caso delle equazioni del regolatore. In particolare, dati $A \in \mathbb{R}^{n \times n}$, $B \in \mathbb{R}^{n \times m}$, $C \in \mathbb{R}^{q \times n}$, $Q \in \mathbb{R}^{q \times r}$, $S \in \mathbb{R}^{r \times r}$, $P \in \mathbb{R}^{n \times r}$, vogliamo calcolare $\Pi \in \mathbb{R}^{n \times r}$ e $\Gamma \in \mathbb{R}^{m \times r}$ tali che:
\begin{equation}
\begin{aligned}
\Pi S &= A\Pi + B\Gamma + P \\
0 &= C\Pi + Q
\end{aligned}
\end{equation}
Come primo passo, riscriviamo la (6) nella forma standard (1). A tal fine, poniamo:
\[
X = \begin{bmatrix} \Pi \\ \Gamma \end{bmatrix} \in \mathbb{R}^{(n+m) \times r}
\]
Questo porta alla seguente riscrittura della (6):
\begin{equation}
\begin{aligned}
\begin{bmatrix} I & 0 \end{bmatrix} XS - \begin{bmatrix} A & B \end{bmatrix} X &= P \\
\begin{bmatrix} C & 0 \end{bmatrix} X &= -Q
\end{aligned}
\end{equation}
che è equivalente a:
\begin{equation}
\underbrace{\begin{bmatrix} -A & -B \\ C & 0 \end{bmatrix}}_{L_1} X + \underbrace{\begin{bmatrix} I & 0 \\ 0 & 0 \end{bmatrix}}_{L_2} XS = \underbrace{\begin{bmatrix} P \\ -Q \end{bmatrix}}_{N}
\end{equation}
dove $L_1 \in \mathbb{R}^{(n+q) \times (n+m)}$, $L_2 \in \mathbb{R}^{(n+q) \times (n+m)}$, e $N \in \mathbb{R}^{(n+q) \times r}$. A questo punto, applicando l'apparato sviluppato in precederza, otteniamo:
\begin{equation}
\left( (S^\top \otimes L_2) + (I_r \otimes L_1) \right) \text{vec}(X) = \text{vec}(N)
\end{equation}


\end{document}



